\section{Introduction}

Hermitian--Yang--Mills flow on a compact Kähler surface can
converge away from finitely many points while its metric
degenerates along a fixed holomorphic filtration. The
limiting connection and point cycle identify the smooth
graded bundle and the lost curvature charge. They do
not, by themselves, determine the rate of the metric
degeneration or the radius on which that charge
concentrates. We determine both for a rank-two
extension with distinct equal-degree line factors
and reduced point defects.

Let \((X,\omega)\) be a connected compact Kähler surface,
and suppose that a locally free rank-two bundle \(V\)
fits into an exact sequence
\begin{equation}\label{eq:intro-extension}
 0\longrightarrow L\xrightarrow{\iota}V
       \longrightarrow\mathcal I_WM\longrightarrow0 .
\end{equation}
The line bundles \(L,M\) have equal degree, and \(W\)
is a nonempty finite scheme. Existence of the locally
free middle bundle is a hypothesis. Fix HE metrics
\(h_L,h_M\), the induced determinant metric \(h_Lh_M\),
and the inclusion \(\iota\). Write
\[
 \dd V=\omega^2/2,\qquad v_0=\int_X\dd V,\qquad
 \kappa=2\pi\deg_\omega L/v_0 .
\]
The metric flow is
\[
 H^{-1}\partial_TH=-2(i\Lambda_\omega F_H-\kappa I),
 \qquad \det H(0)=h_Lh_M .
\]
Its inclusion norm, mass and full derivative energy are
\begin{equation}\label{eq:intro-observables}
 u(T)=\iota^*H(T)/h_L,\qquad
 Z(T)=\int_Xu(T)\,\dd V,\qquad
 \mathcal E(T)=\int_X|D'_{H(T)\otimes h_L^{-1}}\iota|^2\,\dd V .
\end{equation}
In particular \(Z'=-2\mathcal E\).

At a reduced zero \(p\), use holomorphic coordinates
orthonormal for the base metric at \(p\). The norm
of the inclusion in any smooth reference metric
of the fixed determinant has quadratic term \(z^*Q_pz\),
with \(Q_p>0\). The determinant of this form is
independent of that reference metric. Put
\begin{equation}\label{eq:intro-constants}
 \lambda_p=\sqrt{\det Q_p},\qquad
 \sigma=\sum_{p\in W}\lambda_p^{-1},\qquad
 \alpha(T)=Z(T)/v_0,\qquad
 r_p(T)=\sqrt{\alpha(T)/\lambda_p}.
\end{equation}
These constants retain the prescribed inclusion,
line metrics, determinant and base metric.

\begin{theorem}\label{thm:main}
Suppose \(L\not\simeq M\) and \(W\) is reduced.
For every smooth initial metric of determinant
\(h_Lh_M\), the flow exists for all \(T\ge0\) and
has the following properties.
\begin{enumerate}
\item The actual inclusion mass and derivative energy satisfy
\begin{equation}\label{eq:intro-mass}
 Z(T)\sim\frac{v_0^2}{4\pi^2\sigma T},\qquad
 cT^{-2}\le\mathcal E(T)\le CT^{-2}
\end{equation}
for large \(T\). The leading mass coefficient is
independent of the initial metric.
\item There is a fixed smooth reference metric \(H_0\)
of the prescribed determinant, with
\(u_0=\iota^*H_0/h_L\), and constants \(C>0\),
\(0<\gamma<1\), such that
\begin{equation}\label{eq:intro-profile}
 \sup_X\left|\frac{u(T)}{Z(T)}
       -\frac1{v_0}\frac{u_0}{u_0+\alpha(T)}\right|
 \le CT^{-\gamma}.
\end{equation}
For small \(r>0\) and large \(T\), with constants
independent of \(r,T\),
\[
 \sup_{X\setminus B_r(W)}
 \left|\frac{u(T)}{Z(T)}-\frac1{v_0}\right|
 \le C\{T^{-\gamma}+(Tr^2)^{-1}\}.
\]
\item In fixed normalized holomorphic coordinates and
frames near \(p\), the metrics \(H(T,r_p(T)y)\)
converge smoothly on compact \(y\)-sets to
\begin{equation}\label{eq:intro-local-metric}
 H_{\mathrm{loc}}(y)
 =(1+|y|^2)I-\frac{|y|^2+2}{1+|y|^2}YY^*,
 \qquad Y=(y_1,-y_2)^{\mathsf t}.
\end{equation}
This determinant-one metric is anti-self-dual and
has curvature energy \(8\pi^2\). More precisely,
for every fixed \(R>0\),
\[
 \int_{B_{Rr_p(T)}(p)}|F_{H(T)}|^2\,\dd V
 \longrightarrow
 8\pi^2\frac{R^4(R^2+3)}{(1+R^2)^3}.
\]
The first inclusion half-height radius on each
radial ray, defined by \(u(T)=\alpha(T)/2\),
and the first ball radius with curvature energy
\(4\pi^2\), are both \(r_p(T)(1+o(1))\);
the first equivalent is uniform in direction.
\end{enumerate}
\end{theorem}

The local coordinates are Kähler normal at \(p\).
The frames normalize the line weights to one at \(p\)
and put the inclusion in the form
\(\sqrt{\lambda_p}(z_1,-z_2)^{\mathsf t}\);
Lemma~\ref{lem:normalized-reference} constructs them
without changing the determinant.
The theorem follows from
Theorems~\ref{thm:actual-comparison} and
\ref{thm:energy-curvature}.
The proof also accounts for all defect charge in
the ordered limits \eqref{eq:ordered-tightness}.
The energy assertion in \eqref{eq:intro-mass}
is an order estimate, not a differentiated version
of the mass equivalent.

For a general finite defect we obtain a separate
one-sided statement, allowing repeated factors.

\begin{theorem}\label{thm:general-intro}
Under \eqref{eq:intro-extension}, without reducedness
or distinctness of the line factors, every smooth
fixed-determinant flow satisfies \(Z(T)=O(T^{-1})\).
If \(W\) has a nonreduced point, then
\[
 \frac{-Z'(T)}{Z(T)^2}\longrightarrow+\infty,
 \qquad TZ(T)\longrightarrow0.
\]
\end{theorem}

This is Theorem~\ref{thm:general-upper-rate}.
The nonreduced conclusion makes the restriction
in Theorem~\ref{thm:main} mathematically necessary
for its positive mass coefficient. Repeated line
factors present a different issue: the real small
connection eigenspace has dimension three, whereas
the scalar trajectory used here requires dimension one.

\subsection{Relation to earlier work}

Daskalopoulos--Wentworth identify the canonical
connection limit of Yang--Mills flow on Kähler
surfaces \cite{DaskalopoulosWentworth04}, with the
correction \cite{DaskalopoulosWentworth07c}, and identify
the blowup set and its multiplicities
\cite{DaskalopoulosWentworth07b}.
Sibley--Wentworth extend the analytic--algebraic
cycle comparison to compact Kähler manifolds
\cite{SibleyWentworth15}. We use their corrected
multiplicity, which at an isolated surface point
is the full length of the defect sheaf.
These results determine the reflexive graded
connection \(D_{h_L}\oplus D_{h_M}\) and the
point charges in our setting.

The predecessor \cite{W_Singular} treats
\eqref{eq:intro-extension} with arbitrary nonempty
finite \(W\). It proves
\(Z=o(T^{-1/2})\) and fixed-compact convergence
of \(u/Z\) to \(1/v_0\), using a residue estimate
for every splitting. Its remaining rate question
asks whether \(\log Z/\log T\) can still tend
to \(-1/2\), particularly for its toric tangent
bundle. Theorem~\ref{thm:main} gives the limit
\(-1\) in the reduced, distinct-factor class,
including that example. Theorem~\ref{thm:general-intro}
also strengthens the one-sided estimate throughout
the predecessor's class.

Haiden--Katzarkov--Kontsevich--Pandit develop
iterated-weight asymptotics for finite metric
flows \cite{HKKP23} and prove an actual bundle-flow
theorem on compact Riemann surfaces
\cite[Theorem~1.1]{HKKP18}. Their comparison method
for integrably forced metric flows also applies
in higher complex dimension
\cite[Section~4.2]{HKKP18}.
Lam proves bounded-remainder asymptotics on a
fixed compact Kähler base when the graded object
is locally free, under the stated rationality
condition on the induced invariant inner product
\cite[Theorem~1.1, Theorem~8.2]{Lam26}.
The singular-graded extension is discussed in
\cite[Remark~8.3]{Lam26}, rather than included in
that theorem. Our quotient \(\mathcal I_WM\)
is singular, and the coefficient and spatial
profile depend on its local presentation.

The charge-one instanton is classical. We compute
its holomorphic metric, inclusion norm and
curvature in the normalization needed for the
compact extension. For quantitative infinite-time
bubbling on \(\RR^4\), Sire--Wei--Zheng construct
Yang--Mills solutions with selected initial
data and exponentially shrinking scale
\cite[Theorem~1]{SireWeiZheng26}.
Their base, data assertion and scale-selection
problem differ from the fixed compact holomorphic
extension considered here. The radius
\(r_p(T)\asymp T^{-1/2}\) is selected by the
compact inclusion mass and its determinant data.

\subsection{The proof}

The starting point is an explicit smooth family
\(H_a\) on the original compact bundle. It has
inclusion norm \(au_0/(u_0+a)\), an exact
full-derivative identity, and the local limit
\eqref{eq:intro-local-metric}. A level-set residue
estimate gives a lower bound for the actual
dissipation, including the strict nonreduced
conclusion. The compact family itself is not
a solution of the metric flow.

The analytic step is to correct this family
orthogonally to the full small eigenspace of
the real operator \(\frac12D_a^*D_a\) on
Hermitian trace-free endomorphisms. A fixed-base
Sobolev estimate excludes additional modes
concentrated at the points. Heat domination
and the complementary gap give a logarithmic
weighted inverse. The resulting nonlinear
correction has uniformly small mean curvature,
and parameter estimates identify its remaining
one-dimensional curvature with its tangent
up to an integrable error. The small eigenvalue
has leading coefficient \(2\pi^2\sigma/v_0\).

For an arbitrary initial metric, invariant
metric distance first stays bounded by the
forced comparison inequality. The small positive
eigenvalue then gives polynomial decay of that
distance, proving the mass coefficient and the
global profile. Convexity of the Donaldson
functional supplies the actual derivative-energy
upper bound. Finally, rescaled elliptic and
parabolic estimates upgrade the metric comparison
to smooth inner convergence; Chern--Weil theory
and the explicit local curvature integral
identify the concentration radius and defect charge.
